\documentclass[../../main.tex]{subfiles}
\begin{document}
% SOURCE p=260; printed=245; object=ch08-approximate-identity-definition; region=middle
\begin{definition}[逼近单位]\label{definition:ch08-approximate-identity}
前两个定理的大多数应用取 $a=1$，不过 $a=0$ 的情形也有用。当 $\phi\in L^1$、$\int\phi=1$ 时，称 $\{\phi_t\}_{t>0}$ 为一个\textbf{逼近单位}，因为它通过卷积算子给出 $L^p$ 上恒等算子的逼近。这种构造可用于以具有指定正则性的函数逼近 $L^p$ 函数。
\end{definition}
% SUPPLEMENT: id=v14-ch08-approximate-identity-motivation; parent=ch08-approximate-identity-definition; kind=motivation+hypotheses; origin=independent-derivation
\begin{remark}
\textbf{动机与条件（重构新增）。}积分归一化为一保证常数因子不变，尺度收缩保证远离原点的绝对质量趋零，两者合在一起使 $f*\phi_t$ 逼近 $f$。这里指前定理已说明的收敛方式：所有有限指数 $L^p$ 的范数收敛，以及有界一致连续函数的一致收敛；并非对所有 $L^\infty$ 元素都有范数收敛。非负性不在这个定义的必要条件内，但若再要求非负，就有保正及平均解释。$a=0$ 时，同样估计给 $f*\phi_t\to0$；这是积分相消而非保留常数的核，后面也有用途。
\end{remark}
% END SUPPLEMENT: id=v14-ch08-approximate-identity-motivation
% SOURCE p=260; printed=245; object=ch08-smooth-density; number=8.17; region=middle
例如有如下两个重要结果。
\begin{proposition}\label{proposition:ch08-smooth-density}
$C_c^\infty$（从而 $\mathcal S$）在 $L^p$（$1\leqslant p<\infty$）以及 $C_0$ 中稠密，后者取一致范数。
\end{proposition}
\begin{proof}
给定 $f\in L^p$、$\varepsilon>0$，由\cref{proposition:ch07-cc-lp-density}，存在 $g\in C_c$ 使 $\|f-g\|_p<\varepsilon/2$。选 $\phi\in C_c^\infty$ 满足 $\int\phi=1$；例如用\cref{eq:ch08-bump-function}中的 $\psi$，取 $\phi=(\int\psi)^{-1}\psi$。由卷积支撑性质及光滑性命题，$g*\phi_t\in C_c^\infty$；由\cref{theorem:ch08-approximate-identity-convergence}，$t$ 足够小时 $\|g*\phi_t-g\|_p<\varepsilon/2$。故 $\|g*\phi_t-f\|_p<\varepsilon$。

若将 $L^p$ 换成 $C_0$、$\|\cdot\|_p$ 换成 $\|\cdot\|_u$，并用\cref{proposition:ch04-cc-dense-czero}代替 $L^p$ 稠密性结果，同一论证成立。
\end{proof}
% SUPPLEMENT: id=v14-ch08-smooth-density-details; parent=ch08-smooth-density; kind=omitted-proof+motivation+hypotheses; origin=independent-derivation
\begin{remark}
\textbf{归一化、支撑与一致版本的补证（重构新增）。}$\psi$ 非负、连续、紧支撑且在原点严格正，故其积分有限且严格大于零，归一化确实合法。伸缩保持光滑性、非负性及积分一，且 $\phi_t$ 支撑紧。$g\in C_c$ 因而可积，$\phi_t$ 的全部导数有界，光滑卷积命题逐阶给 $C^\infty$。两个紧支撑的和紧，所以卷积也紧支撑。

在 $C_0$ 情形，先取 $g\in C_c$ 使 $\|f-g\|_u<\varepsilon/2$。由前节引理，$g$ 有界且一致连续，因此逼近定理的 (b) 确实适用，给 $\|g*\phi_t-g\|_u<\varepsilon/2$。相加完成一致逼近。又 $C_c^\infty\subseteq\mathcal S$，而 $\mathcal S\subseteq L^p\cap C_0$：$L^p$ 已证，$C_0$ 由零阶的快速衰减得到。因此更大子空间 $\mathcal S$ 也在两个目标空间中稠密。

动机是先截断与连续逼近，再卷积磨光；先做紧支撑逼近保证最后保留紧支撑。$p<\infty$ 用于 $C_c$ 的 $L^p$ 稠密性和卷积近似，不能以此证明 $C_c^\infty$ 在 $L^\infty$ 中稠密；常数函数 $1$ 与任意紧支撑函数的本质一致距离至少为一。
\end{remark}
% END SUPPLEMENT: id=v14-ch08-smooth-density-details
% SOURCE p=260; printed=245; object=ch08-smooth-cutoff; number=8.18; region=middle-lower
\begin{proposition}[$C^\infty$ Urysohn 引理]\label{proposition:ch08-smooth-cutoff}
若 $K\subseteq\mathbb R^n$ 紧，$U$ 为包含 $K$ 的开集，则存在 $f\in C_c^\infty$，使 $0\leqslant f\leqslant1$、$f=1$ 在 $K$ 上成立，且 $\operatorname{supp}f\subseteq U$。
\end{proposition}
\begin{proof}
在 $K\ne\varnothing$、$U^c\ne\varnothing$ 时，令 $\delta=\operatorname{dist}(K,U^c)>0$，并令 $V=\{x:\operatorname{dist}(x,K)<\delta/3\}$。选取非负 $\phi\in C_c^\infty$，使 $\int\phi=1$ 且当 $|x|\geqslant\delta/3$ 时 $\phi(x)=0$；例如取
$\phi=(\int\psi)^{-1}\psi_{\delta/3}$，其中 $\psi$ 如\cref{eq:ch08-bump-function}。置 $f=\chi_V*\phi$。由支撑性质及光滑性，$f\in C_c^\infty$。并且 $0\leqslant f\leqslant1$、$f=1$ 在 $K$ 上成立，且
\begin{align*}
\operatorname{supp}f\subseteq
\{x:\operatorname{dist}(x,K)\leqslant2\delta/3\}\subseteq U.
\end{align*}
退化情形及这些性质的逐项核验如下。
\end{proof}
% SUPPLEMENT: id=v14-ch08-smooth-cutoff-details; parent=ch08-smooth-cutoff; kind=omitted-proof+motivation+hypotheses; origin=independent-derivation
\begin{remark}
\textbf{距离、边界与完整核验（重构新增）。}若 $K=\varnothing$，取 $f=0$；若 $U=\mathbb R^n$ 且 $K\ne\varnothing$，任取有限 $\delta>0$，以下构造同样成立。因此不将到空集的无穷距离直接拿来作伸缩参数。其余情形下，$x\mapsto\operatorname{dist}(x,U^c)$ 为连续正函数，在紧集 $K$ 上达到正最小值，且因 $U^c\ne\varnothing$ 而有限，这证明原证的 $\delta$ 合法。

$V$ 开且有界，故 $\chi_V\in L^1$，其支撑包含于紧集 $\{x:\operatorname{dist}(x,K)\leqslant\delta/3\}$。$\phi$ 有界光滑、全部导数有界，因此卷积光滑。它的支撑包含于上述紧集与闭球 $\overline B(0,\delta/3)$ 的和；对和集中的任一点，三角不等式给到 $K$ 的距离至多 $2\delta/3$，所以支撑紧且包含于 $U$。最后一个包含也可直接反证：若 $x\in U^c$，则到 $K$ 的距离至少为 $\delta$。

非负性与归一化给
\begin{align*}
0\leqslant f(x)=\int\chi_V(x-y)\phi(y)\,\mathrm dy\leqslant\int\phi=1.
\end{align*}
若 $x\in K$，则对每个 $\phi(y)\ne0$ 的点有 $|y|<\delta/3$，故 $\operatorname{dist}(x-y,K)\leqslant|y|<\delta/3$，即 $x-y\in V$。所以被积函数等于 $\phi(y)$，得到 $f(x)=1$；球面上 $\phi=0$，不会留下边界问题。

动机是把一个较小邻域的示性函数磨光，同时预留两层距离，以免支撑越过 $U$ 的边界。紧性提供有界邻域及到闭补集的正距离，开性提供容纳支撑的空间，核的非负性保证 $0\leqslant f\leqslant1$，积分为一保证在 $K$ 上恰等于一，紧支撑控制传播范围。这些条件各负其责，而不是仅凭“卷积光滑”就能推出全部结论。
\end{remark}
% END SUPPLEMENT: id=v14-ch08-smooth-cutoff-details
\end{document}
